\section{ADSD: Attack Algorithm and Implementation}
\label{sec:attack}

\subsection{Single-Context Soft-Prefix Attack}

The first implemented pipeline targets the \emph{first verifier step}, since
early rejection is the most damaging to token-wise speculative throughput. Let
$p^U$ and $q^U$ denote the target and draft next-token distributions after the
relaxed prefix $U$. ADSD optimizes
\begin{align*}
J(U) ={}&
\underbrace{\sum_{y} q^U(y)\,\mathrm{softplus}\!\left(\log q^U(y)-\log p^U(y)\right)}_{\text{soft collapse}}
+ \beta \,\tv(p^U, q^U)
+ \eta \,\mathrm{KL}(q^U \Vert p^U) \\
&- \tau \,\frac{1}{m}\sum_{t=1}^{m} H(u_t)
- \lambda \,\rho(U).
\end{align*}

The soft-collapse term rewards draft overconfidence on tokens that the target
assigns much smaller probability. The TV and reverse-KL terms enlarge the
surrounding disagreement region, the entropy term sharpens the relaxed simplex
vectors, and the projection penalty is the theory-guided term mandated by the
relaxation-gap analysis.

In the implementation, the prefix is parameterized by per-position vocabulary
logits. A softmax converts those logits into simplex vectors, and Adam updates
the logits directly. Each optimization step uses a straight-through argmax
estimator: the forward pass uses the hard projected token at each position,
while the backward pass differentiates through the softmax relaxation. The
released code therefore optimizes over the full vocabulary while keeping the
forward path faithful to the discrete token manifold used at deployment time.

Although the theory is phrased in terms of explicit nearest-neighbor
projection, the implementation uses argmax projection for computational
efficiency. The theoretical link is preserved by retaining the explicit
$L_2$ projection-distortion penalty $\rho(U)$ in the objective. Model selection
is verifier-aligned: although the optimizer follows the smooth objective
$J(U)$, the retained attack is ranked by its exact projected first-token
acceptance
\[
\alpha_1 = \min\!\left(1, \frac{p(y_1)}{q(y_1)}\right),
\]
computed after hard projection in the actual speculative-decoding pipeline.

In the experiments reported here, the core hyperparameters are a five-token
prefix, Adam learning rate $0.3$, temperature $1.0$, entropy weight $0.01$, TV
weight $1.0$, reverse-KL weight $0.1$, and gradient clipping at $1.0$. The
projection weight $\lambda$ is calibrated automatically from the ratio between
the unpenalized objective and the initial projection distortion, then swept on
Palmetto through a multiplicative factor. Because the search landscape is
highly non-convex, the implementation warm-starts the initial logits from a
short manual adversarial seed and searches the full vocabulary around that
initialization.

\subsection{Calibration-Batch Universal Prefix Extension}

To move beyond a single boundary context, ADSD also optimizes a shared
universal prefix over a small calibration batch of GSM8K training prompts. For
each calibration prompt $x^{(j)}$, the target model first produces a fixed
continuation $y^{*(j)}_{1:K}$ of length $K$. These target-only rollouts define
the temporal contexts on which the universal prefix is optimized.

Let $p_i^{U,j}$ and $q_i^{U,j}$ denote the target and draft next-token
distributions at timestamp $i$ on calibration prompt $j$. The universal-prefix
objective is
\begin{align*}
J_{\mathrm{UAT}}(U) ={}&
\frac{1}{N}\sum_{j=1}^{N}
\Bigg[
\lambda_c \sum_{i=1}^{s} \bar{w}_i\,
\mathrm{SC}\!\left(q_i^{U,j}, p_i^{U,j}\right)
+ \beta \sum_{i=1}^{K} w_i\,\tv\!\left(q_i^{U,j}, p_i^{U,j}\right)
\\
&\qquad\qquad
+ \eta \sum_{i=1}^{K} w_i\,\mathrm{KL}\!\left(q_i^{U,j} \Vert p_i^{U,j}\right)
\Bigg]
- \tau \,\frac{1}{m}\sum_{t=1}^{m} H(u_t)
- \lambda \,\rho(U),
\end{align*}
where $w_i \propto K-i+1$ linearly emphasizes earlier timestamps and
$\bar{w}_i$ is the corresponding normalized restriction to the first $s$
timestamps. The term $\mathrm{SC}$ is the same soft-collapse signal used in the
single-context attack. The design thus preserves the theoretical emphasis on
early verifier failure while extending disagreement pressure across the full
target trajectory.

\subsection{Gradient-Guided Discrete Seed Search}

The universal-prefix implementation also includes a discrete initialization
stage before continuous refinement. Starting from the warm-start seed, ADSD
computes gradient-guided token replacements at each prefix position, retains
the top vocabulary candidates, and runs a small beam search over the prefix.
Candidate prefixes are ranked by exact projected first-token acceptance, with
weighted trajectory TV used as the tie-breaker. The best discrete candidate
then initializes the same straight-through optimization described above.

This discrete seed-search stage is important in practice. It places the
continuous optimizer in a stronger basin, reduces dependence on a single
hand-written seed, and materially improves the quality of the projected
universal prefix recovered after refinement. Empirically, it is this hybrid
discrete-plus-continuous pipeline that produces the strongest universal attacks
reported in the experiments.
